%%
%% Chapter 3: Methodology (方法设计)
%%

\chapter{方法设计}

\section{理论框架}

本研究基于相关理论基础，构建了系统的研究框架。理论框架是研究的基础，它为研究假设的提出和实证分析提供了理论支撑。

\section{研究假设}

根据理论分析和文献综述的结果，本研究提出以下研究假设：

假设1：变量A对变量B具有显著的正向影响。

假设2：变量C在变量A与变量B的关系中起中介作用。

假设3：变量D对变量A与变量B的关系具有调节效应。

\section{研究模型}

基于以上理论分析和研究假设，本研究构建了如下研究模型。该模型包含了自变量、因变量、中介变量和调节变量，能够全面反映各变量之间的关系。
